% Summary of diercks2013: Indirect agree in {Lubukusu} complementizer agreement

\begin{abstract}
\begin{enumerate}
  \item In Lubukusu, complementizers sometimes agree with the most local superordinate subject.
  \item Proposal: this is actually local agreement between the complementizer and a null, subject-oriented anaphor in specCP, bound by the next highest subject as a Long-distance Reflexive.
\end{enumerate}

Thoughts:

\begin{itemize}
  \item Look back at discussion (S4.1.1, 4.1.2) of the semantics of when you use the agreeing vs. non-agreeing complementizer.
  \item Semantics-y long-distance things seem to go Upwards, whereas other long-distance things seem to go Downwards.
\end{itemize}

\begin{enumerate}[start=4]
  \item In Tsez, verbs agree with absolutive arguments in noun class.
  \item When the absolutive argument is a clause, and that clause has an absolutive topic (covertly moved at LF), verbs agree with the absolutive topic. If there is no absolutive topic, they agree with the clause itself.
\end{enumerate}
\end{abstract}

\section*{Data}

\begin{enumerate}
  \item mother [boy bread.III.ABS ate].IV \textbf{IV}-know (PLA)
  \item mother [boy bread.III.ABS ate].IV \textbf{III}-know (LDA)
\end{enumerate}
`The mother knows the boy ate bread.'
